\documentclass[12pt,a4paper]{amsart}

\usepackage{amsmath,amssymb,amsthm}
\usepackage{mathtools}
\usepackage[colorlinks=true,linkcolor=blue!60!black,citecolor=green!50!black,
            urlcolor=blue!70!black]{hyperref}
\usepackage[shortlabels]{enumitem}
\usepackage{booktabs}

%% ---- Theorem environments ----------------------------------------
\newtheorem{theorem}{Theorem}[section]
\newtheorem{lemma}[theorem]{Lemma}
\newtheorem{proposition}[theorem]{Proposition}
\newtheorem{corollary}[theorem]{Corollary}
\theoremstyle{definition}
\newtheorem{definition}[theorem]{Definition}
\newtheorem{remark}[theorem]{Remark}
\newtheorem{assumption}[theorem]{Assumption}
\newtheorem{problem}{Problem}[section]

%% ---- Notation ----------------------------------------------------
\newcommand{\PW}{\mathrm{PW}_\omega}
\newcommand{\Kop}{\mathcal{K}_c}
\newcommand{\GVM}{G^{(N)}_{\mathbf{p},c}}
\newcommand{\QM}{Q^{(N)}_{\mathbf{p},c}}
\newcommand{\psin}{\psi_n^{(c)}}
\newcommand{\EN}{E_{N,c}}
\newcommand{\norm}[1]{\|#1\|}
\newcommand{\ip}[2]{\langle #1,\, #2\rangle}
\newcommand{\Rbulk}{\mathcal{R}_{\mathrm{bulk}}}
\newcommand{\Rturning}{\mathcal{R}_{\mathrm{turn}}}

\title[Frame Coercivity under First-Order Quadrature Control]
      {Frame Coercivity of Discrete Prolate Sampling Operators\\
       under First-Order Quadrature Control:\\
       Defect Decomposition and a Bulk--Turning Decay Mechanism}

\thanks{This version replaces the earlier conditional decay conjecture
by a proof architecture based on a quadrature--projection decomposition,
a fully quantified bulk phase-gap lemma, bulk-localised
integration-by-parts control, and a conservative turning-point estimate
via Cauchy--Schwarz and Airy-scale widths. The remaining quadrature
input is reduced to a standard first-order quadrature consistency
assumption. A companion paper (Paper~II) develops scaling limits, trace
formula, and Weil-operator connections; see~\cite{PaperII}. A further
companion (Paper~II\textsubscript{quad}) develops a PSWF-Gauss/XRY
regime and a conditional implication framework for oscillation-compatible
quadrature; see~\cite{PaperIIquad}. The semiclassical WKB amplitude and
spectral density results used in Section~\ref{sec:decay} are established
unconditionally in Paper~IV~\cite{PaperIV}.}

\author{Anonymous Author}
\address{}
\email{}
\date{May 2026}

\subjclass[2020]{42C05, 42B10, 47B32, 41A60}

\keywords{prolate spheroidal wave functions, finite frames,
sampling in reproducing kernel Hilbert spaces,
Gram operator, Marcinkiewicz--Zygmund inequalities,
frame perturbation, banded decay, coercivity,
WKB asymptotics, turning points, quadrature error}

\begin{document}
\maketitle

\begin{abstract}
Let $\PW$ be the Paley--Wiener space of $\omega$-bandlimited functions
and $V_{N,c}:=\operatorname{span}\{\psi_0^{(c)},\dots,\psi_{N-1}^{(c)}\}$
the $N$-dimensional prolate subspace at time-bandwidth product $c=\omega T$.
We study the discrete sampling form
$\QM(f)=\sum_{j=1}^N w_j|f(p_j)|^2$ as a finite frame form on $V_{N,c}$.

We establish coercivity results for random sampling configurations and
for weighted sampling models satisfying a Marcinkiewicz--Zygmund lower
frame bound. Our main structural result is an exact defect decomposition
reducing frame analysis to simultaneous control of a weighted-frame
defect and a quadrature defect. Assuming a standard first-order
quadrature consistency estimate, we derive
\[
  |E_{mn}| \le C\,hc
  \qquad \text{for all } 0\le m,n\le N-1.
\]
The exact identity $E_{mn} = R_{mn}^{\mathrm{quad}}$
(Proposition~\ref{prop:xry_split}) shows the defect is purely a
quadrature error. Separately, the continuous PSWF product integral
exhibits algebraic decay via a phase-gap estimate and iterated
integration by parts in the bulk sector
$m,n\le\gamma_0 N_{\mathrm{Sh}}$.

The bulk--turning analysis is an independent structural component
identifying the oscillatory decay mechanism; the proof of the main
stability theorem does not depend on it. The transfer of continuous
decay to the discrete level requires the DSTP interface axiom,
verified for concrete quadrature schemes in a companion paper.
\end{abstract}

\tableofcontents

%% ==================================================================
\section{Introduction}
%% ==================================================================

The prolate spheroidal wave functions (PSWFs) form an orthonormal basis
for the finite-dimensional subspace $V_{N,c}$ of the Paley--Wiener
space $\PW$ that is optimally concentrated in a time--frequency window
\cite{Slepian1978,BonamiKaroui}. The weighted evaluation map
$f\mapsto(\sqrt{w_j}f(p_j))_j$ defines a finite frame operator on
$V_{N,c}$, whose associated Gram matrix encodes the stability of the
sampling scheme.

This paper addresses two questions. First, which classes of sampling
configurations yield coercive Gram operators on $V_{N,c}$? Second, how
does the discrete sampling Gram matrix compare to the reference spectral
matrix $G^{\mathrm{ref}}_{N,c}=\Kop|_{V_{N,c}}$ in the PSWF
eigenbasis? For the first question we prove coercivity for random
sampling and for weighted MZ frames. For the second, we establish an
exact algebraic defect decomposition, derive a global first-order
quadrature bound, and separately analyse the oscillatory decay mechanism
via a WKB-bulk/Airy-turning decomposition in the bulk sector.

The central contribution is a \emph{reduction principle}: discrete
frame stability is reduced to a single verifiable spectral transfer
property (DSTP), and the precise obstruction is isolated and made
explicit.

\subsection*{Main results}
The proof architecture has four components:
\begin{enumerate}
\item \textbf{Quadrature--projection split}
  (Proposition~\ref{prop:xry_split}): $E_{mn}$ is a pure quadrature
  remainder.
\item \textbf{Spectral parameter control and bulk bridge}
  (Proposition~\ref{prop:spectral_gap},
  Lemma~\ref{lem:discrete_continuous_bridge}): the two-term asymptotics
  of $\chi_n(c)$ yield explicit spectral gaps; the bulk bridge lemma
  establishes the uniform discrete--continuous correspondence required
  for the phase-gap and IBP arguments.
\item \textbf{Uniform phase gap and IBP}
  (Lemmas~\ref{lem:bulk_spectral_sep}--\ref{lem:bulk_decay_continuous}):
  phase separation, amplitude control, and iterated integration by parts
  yield algebraic decay for the continuous PSWF product integral in the
  bulk sector.
\item \textbf{Quadrature reduction}
  (Assumption~\ref{ass:quad_first_order},
  Theorem~\ref{thm:banded_decay}): first-order quadrature consistency
  yields $|E_{mn}|\le Chc$.
\end{enumerate}

%% FIX 5: DSTP necessity — abgeschwächt auf "within this proof architecture"
\begin{remark}[Necessity of DSTP]\label{rem:dstp_necessity}
DSTP is necessary \emph{within this proof architecture}: any
Schur-summable decay estimate of the form
$|E_{mn}|\le Chc(1+|m-n|)^{-p}$ must encode a condition of DSTP
type, since the structure-blind bound (Theorem~\ref{thm:banded_decay})
is sharp for generic rules (Remark~\ref{rem:ax5_independence}).
Whether every stable scheme must satisfy DSTP or an equivalent
condition is an open question; see Section~\ref{sec:open}.
DSTP is a non-uniform spectral localisation condition in the PSWF
eigenbasis, not a sampling density condition.
\end{remark}

\subsection*{Organisation}
Section~\ref{sec:setup} fixes notation.
Section~\ref{sec:coercivity} proves coercivity.
Section~\ref{sec:decay} establishes the decay mechanism.
Section~\ref{sec:stability} derives frame stability.
Section~\ref{sec:bridge} explains the conceptual boundary.
Section~\ref{sec:open} lists open problems.
Appendix~\ref{app:airy} contains the Airy approximation lemma.

%% ==================================================================
\section{Setup and notation}\label{sec:setup}
%% ==================================================================

\subsection{Paley--Wiener space and concentration operator}

We use $\widehat{f}(\xi)=\int_{\mathbb{R}}f(x)e^{-2\pi ix\xi}\,dx$.
The Paley--Wiener space $\PW$ consists of all $f\in L^2(\mathbb{R})$
with Fourier support in $[-\omega/(2\pi),\omega/(2\pi)]$, with kernel
\[
  K_{\PW}(x,y)
  =\tfrac{\omega}{\pi}\operatorname{sinc}(\omega(x-y)/\pi)
\]
\cite{BonamiKaroui}. The prolate concentration operator is
\[
  \Kop f(x)=\int_{-T}^{T}\frac{\sin\omega(x-y)}{\pi(x-y)}f(y)\,dy,
\]
with eigenvalues $1>\lambda_0(c)>\lambda_1(c)>\cdots>0$ and PSWF
eigenfunctions $\{\psin\}_{n\ge 0}$ \cite{Slepian1978}. We set
$c=\omega T$, $V_{N,c}:=\operatorname{span}\{\psi_0^{(c)},\dots,
\psi_{N-1}^{(c)}\}$, Shannon number $N_{\mathrm{Sh}}:=2c/\pi$, and
fix $\gamma_0\in(0,\tfrac{1}{4})$. The \emph{bulk sector} is
$\mathcal{B}:=\{0\le k\le\gamma_0 N_{\mathrm{Sh}}\}$.

Throughout this paper, $\gamma_0\in(0,\tfrac{1}{4})$ is fixed
independently of $c$ and $N$. All constants depending on $\gamma_0$
are uniform as $c\to\infty$ with $\gamma_0$ held fixed.
The scaling $N\sim c$ means $N/N_{\mathrm{Sh}}\to\pi/2$; the bulk
sector then satisfies $|\mathcal{B}|=\lfloor\gamma_0
N_{\mathrm{Sh}}\rfloor\sim 2\gamma_0 c/\pi$, growing linearly in $c$
with fixed proportion $\gamma_0$.

\subsection{Frame operator and defect matrix}

\begin{definition}[Weighted Gram matrix]\label{def:GramForm}
For $\mathbf{p}=(p_1,\dots,p_N)\subset[-T,T]$ and weights $w_j>0$,
\[
  G^{\mathrm{w}}_{N,c}(\mathbf{p},D)
  :=M(\mathbf{p},c)^*D\,M(\mathbf{p},c),
  \quad M_{jn}:=\psi_n^{(c)}(p_j),\quad
  D:=\operatorname{diag}(w_j).
\]
\end{definition}

\begin{definition}[Reference matrix and defect]\label{def:ref_matrix}
\[
  G^{\mathrm{ref}}_{N,c}
  :=\operatorname{diag}(\lambda_0(c),\dots,\lambda_{N-1}(c)),
  \qquad
  \EN:=G^{\mathrm{w}}_{N,c}(\mathbf{p},D)-G^{\mathrm{ref}}_{N,c}.
\]
\end{definition}

\begin{definition}[Frame-stable regime]\label{def:framestable}
$\mathbf{p}$ is frame-stable if $\norm{\EN}_2<\lambda_{N-1}(c)$.
\end{definition}

%% ------------------------------------------------------------------
\subsection{Spectral asymptotics and the bulk bridge lemma}
%% ------------------------------------------------------------------

\begin{proposition}[Two-term spectral asymptotics and uniform gap]%
\label{prop:spectral_gap}
For $n\in\mathcal{B}$:
\[
  \chi_n(c) = n(n+1) + \frac{(2n+1)c}{\pi}
              + O\!\left(\frac{c^2}{n}\right),
\]
\cite[Section~2.5]{OsipovRokhlinXiao}, \cite[Introduction]{PaperIV}.
For $m\ne n$ in $\mathcal{B}$, subtracting:
\begin{equation}\label{eq:chi_diff}
  \chi_m(c)-\chi_n(c)
  =(m-n)(m+n+1)+\frac{2(m-n)c}{\pi}
   +O\!\left(\frac{c^2}{\min(m,n)}\right).
\end{equation}
We split into two regimes.

\smallskip
\noindent\textbf{Regime~(A): deep bulk, $\min(m,n)\ge c^{1/2}$.}
Here $m+n+1\le 4\gamma_0 c/\pi+1$ and the remainder in
\eqref{eq:chi_diff} satisfies
$c^2/\min(m,n)\le c^{3/2}=o(c|m-n|)$ for $|m-n|\ge 1$.
The linear term $2c(m-n)/\pi$ dominates for $\gamma_0<\tfrac{1}{4}$:
\[
  |\chi_m-\chi_n|
  \ge\frac{c}{\pi}|m-n|
     \underbrace{(2-4\gamma_0-o(1))}_{\ge 1\text{ for }c\ge c_0}
  \ge\frac{c}{\pi}|m-n|.
\]
Hence in Regime~(A):
\[
  |\chi_m-\chi_n|\ge\kappa_1^{(A)}\,c\,|m-n|,
  \qquad\kappa_1^{(A)}:=\frac{1}{\pi}.
\]

\noindent\textbf{Regime~(B): near base, $\min(m,n)<c^{1/2}$.}
Here $m,n\le c^{1/2}$, so $m+n+1\le 2c^{1/2}+1$. The remainder
$c^2/\min(m,n)>c^{3/2}$ can dominate the linear-in-$c$ term, so
the Regime~(A) bound fails. However, the quadratic term alone gives
\[
  |\chi_m-\chi_n|\ge|m-n|,
\]
and therefore:
\[
  |\chi_m-\chi_n|\ge\kappa_1^{(B)}\,|m-n|,
  \qquad\kappa_1^{(B)}:=1.
\]

\noindent\textbf{Uniform constant.}
Setting
\[
  \kappa_1:=\min\!\bigl(\kappa_1^{(A)},\,\kappa_1^{(B)}\bigr)
           =\min\!\Bigl(\tfrac{1}{\pi},\,1\Bigr)
           =\tfrac{1}{\pi},
\]
both regimes satisfy
\begin{equation}\label{eq:chi_gap}
  |\chi_m(c)-\chi_n(c)|\ge\kappa_1\max(1,c)\,|m-n|
  \qquad\forall\,m\ne n\in\mathcal{B},\;c\ge c_0,
\end{equation}
since Regime~(A) gives $\kappa_1^{(A)}c|m-n|\ge\kappa_1 c|m-n|$
and Regime~(B) gives $\kappa_1^{(B)}|m-n|\ge\kappa_1|m-n|$,
and $\max(1,c)$ absorbs both cases uniformly.
The IBP argument in Lemma~\ref{lem:ibp_uniform} depends only on
\emph{positivity} of $\kappa_1$, not on its precise value; the
constant $C_p$ in that lemma is therefore the same for both regimes.
No pointwise Lipschitz condition on $n\mapsto\chi_n$ is used;
\eqref{eq:chi_gap} is a direct algebraic consequence of the two-term
expansion.

%% FIX 2: Regime B deliberate non-sharpness declaration
The bound in Regime~(B) is deliberately non-sharp: for $n\ll c$ one
has $\chi_n\sim n(n+1)$, giving a stronger gap, but the coarser bound
$|\chi_m-\chi_n|\ge|m-n|$ suffices for all subsequent arguments and
avoids case analysis at the spectral base.
There exists $c_0>0$ such that for all $c\ge c_0$ the bound
\eqref{eq:chi_gap} holds uniformly across both regimes, eliminating
any boundary issue at small $c$.
\end{proposition}

\begin{lemma}[Uniform discrete--continuous bulk bridge]%
\label{lem:discrete_continuous_bridge}
Fix $\gamma_0\in(0,\tfrac{1}{4})$ and $\delta\in(0,1)$. For all
$n\in\mathcal{B}$ and $x\in I_\delta:=[-(1-\delta)T,(1-\delta)T]$:
\begin{enumerate}[(i)]
\item \textup{(Classical momentum scale)}
  $\omega_n^{\mathrm{cl}}(x)
  :=\sqrt{\chi_n(c)/T^2-c^2x^2/T^2+c^2}\asymp n$,
  uniformly in $n\in\mathcal{B}$, $c\ge c_0$, $x\in I_\delta$.
\item \textup{(Spectral-to-spatial correspondence)}
  $|\chi_m(c)-\chi_n(c)|
  =T^2|(\omega_m^{\mathrm{cl}})^2(x)-(\omega_n^{\mathrm{cl}})^2(x)|$
  for every $x\in[-T,T]$.
\item \textup{(Phase derivative lower bound)}
  For $x\in\Rbulk^{(m,n)}$ (Definition~\ref{def:bulk}),
  \[
    |\Phi_{mn}'(x)|
    =\frac{|\chi_m(c)-\chi_n(c)|}{T(\sqrt\mu+\sqrt\nu)}
    \ge\frac{\kappa_1\max(1,c)|m-n|}{T\cdot C_\nu\,c}
    \ge\frac{\kappa}{T}|m-n|,
  \]
  where $\mu:=\chi_m-c^2x^2/T^2$, $\nu:=\chi_n-c^2x^2/T^2$, and
  $\sqrt\mu+\sqrt\nu\le C_\nu c$ on $\mathcal{B}$ (since
  $\chi_k\le C\gamma_0^2 c^2$ on $\mathcal{B}$ gives
  $\sqrt{\chi_k}\le C\gamma_0 c$). The factor $c$ in $\max(1,c)$
  cancels with the $c$ in the denominator $C_\nu c$; the resulting
  $\kappa=\kappa_1/C_\nu$ is independent of $c$ and of which
  regime $(m,n)$ falls into.
\item \textup{(Amplitude derivative control, c-uniform)}
  The Pr\"ufer amplitude satisfies $r_n(x)=O(1)$ uniformly,
  and its $k$-th derivative obeys
  $\|\partial_x^k r_n\|_{L^\infty(I_\delta)}\le C_k$
  with $C_k$ independent of $n\in\mathcal{B}$ and $c\ge c_0$,
  by \cite[Lemma~2.4, Lemma~3.5]{PaperIV}.
\item \textup{(Discrete--continuous scaling bridge)}
  The map $n\mapsto\chi_n(c)$ satisfies
  \[
    \Delta\chi_n:=\chi_{n+1}(c)-\chi_n(c)
    =2n+2+\frac{2c}{\pi}
      +O\!\left(\frac{c^2}{n^2}\right),
  \]
  which is the \emph{exact discrete analogue} of the continuous
  derivative $\partial_n\chi_n(c)$. This identity links spectral
  index increments to WKB momentum increments and justifies the
  identification of discrete-index differences with continuous
  phase-gap bounds: no implicit $d/dn\leftrightarrow(1/c)\,d/dx$
  substitution is made anywhere in this paper. All spectral-gap
  bounds are derived purely from the algebraic identity
  \eqref{eq:chi_diff}.
\end{enumerate}
\end{lemma}
\begin{proof}
(i) follows from \cite[Lemma~2.3(i)]{PaperIV} with the identification
$\omega_n^{\mathrm{cl}}(x)^2=c^2(1-x^2/T^2)+\chi_n/T^2\ge\chi_n/T^2
\ge C_\gamma n^2/T^2$ for $n\in\mathcal{B}$ (using $\chi_n\ge Cn^2$).
(ii) is immediate from Definition~\ref{def:bulk}.
(iii): numerator from \eqref{eq:chi_gap}; denominator bounded by
$\sqrt\mu+\sqrt\nu\le 2\max(\sqrt{\chi_m},\sqrt{\chi_n})\le C\gamma_0 c
\le C_\nu c$ since $m,n\in\mathcal{B}$ implies $\chi_k\le C_{\gamma_0}c^2$.
(iv): see \cite[Lemma~2.4, Lemma~3.5]{PaperIV}, which establishes
uniform $C^k$ control of the Pr\"ufer amplitude on $I_\delta$, with all
constants depending on $\delta$, $\gamma_0$, $k$ but not on $c$.
(v): subtract consecutive instances of the two-term formula in
Proposition~\ref{prop:spectral_gap}.
\end{proof}

\begin{remark}[Role of the bridge lemma]
Lemma~\ref{lem:discrete_continuous_bridge} is the structural backbone
connecting spectral index geometry (discrete) to WKB phase geometry
(continuous). Part~(iii) makes explicit that $c$ cancels in the
phase-gap bound --- not by assumption, but by the simultaneous
$c$-scaling of numerator ($\kappa_1 c|m-n|$) and denominator
($C_\nu c$); this cancellation is uniform across both Regime~(A)
and Regime~(B) of Proposition~\ref{prop:spectral_gap}.
Part~(iv) prevents spurious $c^{k/2}$ growth in the IBP
amplitude chain. Part~(v) eliminates the need for any
Bohr--Sommerfeld difference argument or implicit differentiability
of $n\mapsto\chi_n$. Every subsequent lemma in Section~\ref{sec:decay}
cites exactly one part of this lemma and nothing else for its
scaling inputs.
\end{remark}

%% ==================================================================
\section{Unisolvence and coercivity}\label{sec:coercivity}
%% ==================================================================

\begin{definition}[Frame condition]\label{def:frame}
$\{(p_j,w_j)\}_{j=1}^N$ is a weighted frame for $V_{N,c}$ if
$A\norm{f}^2\le\sum_j w_j|f(p_j)|^2\le B\norm{f}^2$,
$0<A\le B<\infty$, for all $f\in V_{N,c}$.
\end{definition}

\begin{definition}[Unisolvent set]
$\mathbf{p}$ is unisolvent for $V_{N,c}$ if
$f\mapsto(f(p_1),\dots,f(p_N))$ is injective on $V_{N,c}$.
\end{definition}

\begin{lemma}[Identity principle]\label{lem:identity}
Every $f\in\PW$ is entire of exponential type; if it vanishes on a set
with a limit point, then $f\equiv 0$.
\end{lemma}
\begin{proof}
$\PW$ embeds into the Bernstein space $B_\omega^\infty$ of entire
functions of exponential type $\omega$ \cite[Ch.~6]{BonamiKaroui};
the claim is the standard identity theorem for such functions.
\end{proof}

\begin{proposition}[MZ lower bound implies unisolvence]%
\label{prop:mz_unisolvent}
A weighted frame with lower bound $A>0$ is unisolvent.
\end{proposition>
\begin{proof}
Suppose $f\in V_{N,c}$ satisfies $f(p_j)=0$ for all $j$.
Then $\sum_j w_j|f(p_j)|^2=0$, so the lower frame bound gives
$A\norm{f}^2\le 0$, hence $f=0$. Thus the evaluation map is injective.
\end{proof}

\begin{proposition}[Random sampling]\label{prop:random}
Let $p_1,\dots,p_N$ be i.i.d.\ absolutely continuous on $[-T,T]$.
Then $\mathbf{p}$ is unisolvent for $V_{N,c}$ almost surely.
\end{proposition>
\begin{proof}
$\det M(\mathbf{p},c)$ is a real-analytic function of $(p_1,\dots,p_N)$
and is non-identically zero (it is non-zero for any Chebyshev-type
configuration). Since the distribution of $\mathbf{p}$ is absolutely
continuous, $\det M(\mathbf{p},c)=0$ has probability zero.
\end{proof}

\begin{lemma}[Pointwise RKHS bound]\label{lem:upper}
$|f(x)|^2\le(\omega/\pi)\norm{f}_{L^2}^2$ for all $f\in\PW$.
\end{lemma}
\begin{proof}
By the reproducing kernel property,
$f(x)=\langle f,K_{\PW}(\cdot,x)\rangle$ and
$\norm{K_{\PW}(\cdot,x)}_{L^2}^2=K_{\PW}(x,x)=\omega/\pi$.
Cauchy--Schwarz gives $|f(x)|^2\le(\omega/\pi)\norm{f}_{L^2}^2$.
\end{proof}

\begin{lemma}[Operator norm bound]\label{lem:lambda_max_bound}
$\lambda_{\max}(G^{\mathrm{w}}_{N,c})\le N\omega/\pi$.
\end{lemma}
\begin{proof}
For any $\mathbf{a}\in\mathbb{C}^N$ with $\norm{\mathbf{a}}=1$,
set $f=\sum_n a_n\psi_n^{(c)}\in V_{N,c}$. Then
\[
  \mathbf{a}^*G^{\mathrm{w}}_{N,c}\mathbf{a}
  =\sum_j w_j|f(p_j)|^2
  \le\frac{\omega}{\pi}\norm{f}_{L^2}^2\sum_j w_j
  \le\frac{\omega}{\pi}\cdot N\cdot\frac{2T}{N}
  \cdot\norm{f}_{L^2}^2,
\]
using Lemma~\ref{lem:upper} pointwise and $\sum_j w_j=2T$,
$\norm{f}_{L^2}^2=1$. This gives $\lambda_{\max}\le N\omega/\pi$.
\end{proof}

\begin{theorem}[Frame coercivity]\label{thm:coercivity}
If $\mathbf{p}$ is unisolvent and $D=I$,
\[
  \sqrt{\lambda_{\min}(\GVM)}\,\norm{f}_{L^2}
  \le\QM(f)^{1/2}
  \le\sqrt{N\omega/\pi}\,\norm{f}_{L^2}
  \qquad\forall f\in V_{N,c}.
\]
\end{theorem}
\begin{proof}
The upper bound follows from Lemma~\ref{lem:lambda_max_bound}.
For the lower bound: since $\mathbf{p}$ is unisolvent,
$M(\mathbf{p},c)$ has full column rank $N$, so
$G^{\mathrm{w}}_{N,c}=M^*M$ is positive definite with
$\lambda_{\min}(\GVM)>0$. Then
$\QM(f)=\mathbf{a}^*G^{\mathrm{w}}_{N,c}\mathbf{a}
\ge\lambda_{\min}(\GVM)\norm{\mathbf{a}}^2
=\lambda_{\min}(\GVM)\norm{f}_{L^2}^2$,
where $\mathbf{a}$ is the coefficient vector of $f$ in the PSWF basis.
\end{proof}

%% ==================================================================
\section{Quadrature decomposition and oscillatory decay}\label{sec:decay}
%% ==================================================================

\begin{proposition}[Quadrature--Projection Decomposition]%
\label{prop:xry_split}
\[
  E_{mn}
  =\underbrace{\Bigl(
     \int_{-T}^{T}\psi_m^{(c)}\overline{\psi_n^{(c)}}\,dx
     -\lambda_n(c)\delta_{mn}\Bigr)}_{=\,0}
   +R_{mn}^{\mathrm{quad}},
\]
where $R_{mn}^{\mathrm{quad}}:=\sum_j w_j f_{mn}(p_j)
-\int_{-T}^T f_{mn}\,dx$ and
$f_{mn}:=\psi_m^{(c)}\overline{\psi_n^{(c)}}$.
%% FIX 1: explicit eigenrelation step
Since $\{\psi_n^{(c)}\}$ are orthonormal in $L^2(\mathbb{R})$ and
satisfy $\Kop\psi_n^{(c)}=\lambda_n(c)\psi_n^{(c)}$, taking the
$L^2(\mathbb{R})$-inner product with $\psi_m^{(c)}$ gives
\[
  \int_{-T}^{T}\psi_m^{(c)}\overline{\psi_n^{(c)}}\,dx
  =\langle\psi_m^{(c)},\Kop\psi_n^{(c)}\rangle_{L^2(\mathbb{R})}
  =\lambda_n(c)\langle\psi_m^{(c)},\psi_n^{(c)}\rangle_{L^2(\mathbb{R})}
  =\lambda_n(c)\delta_{mn},
\]
so the bracketed term vanishes exactly, and
$E_{mn}=R_{mn}^{\mathrm{quad}}$ is a pure quadrature error.
\end{proposition}

\begin{definition}[Bulk and turning regions]\label{def:bulk}
Let $x_n^*:=(T/c)\sqrt{\chi_n(c)}$ be the positive turning point of
$\psi_n^{(c)}$ and fix $\rho>0$. Set
\[
  U_n:=\{|x\mp x_n^*|\le\rho c^{-1/3}\}\cap[-T,T],
  \quad
  \Rturning^{(m,n)}:=U_m\cup U_n,
  \quad
  \Rbulk^{(m,n)}:=[-T,T]\setminus\Rturning^{(m,n)}.
\]
\end{definition}

\begin{remark}
The Airy width $O(c^{-1/3})$ is the standard transition scale
\cite{Olver1974}. The turning point $x_n^*$ is well-defined for
$n\in\mathcal{B}$ since $\chi_n(c)<c^2$ there
\cite[Section~2.3]{OsipovRokhlinXiao}.
\end{remark}

\begin{lemma}[Bulk spectral separation]\label{lem:bulk_spectral_sep}
For $x\in\Rbulk^{(m,n)}$ and $k\in\{m,n\}$:
\[
  \left|\chi_k(c)-\frac{c^2x^2}{T^2}\right|\ge\delta_0 c^{2/3},
  \qquad\delta_0:=\rho^2/T^2.
\]
\end{lemma>
\begin{proof}
$\chi_k-c^2x^2/T^2=(c^2/T^2)(x_k^*-x)(x_k^*+x)$; on
$\Rbulk^{(m,n)}$ both factors $\ge\rho c^{-1/3}$ by
Definition~\ref{def:bulk}.
\end{proof}

\begin{lemma}[Uniform phase gap]\label{lem:phase_gap_final}
For $m\ne n$ in $\mathcal{B}$ and $x\in\Rbulk^{(m,n)}$:
\[
  |\Phi_{mn}'(x)|\ge\frac{\kappa}{T}|m-n|,
\]
with $\kappa>0$ depending only on $\gamma_0$, $T$, independent of
$c\ge c_0$ and $m,n\in\mathcal{B}$.
\end{lemma}
\begin{proof}
Write $\mu=\chi_m-c^2x^2/T^2$, $\nu=\chi_n-c^2x^2/T^2$.
By Lemma~\ref{lem:bulk_spectral_sep}, $\sqrt\mu+\sqrt\nu\le
2\sqrt{\chi_{N-1}}\le C_\nu c$.

\noindent\textbf{Numerator.}
By Proposition~\ref{prop:spectral_gap} and \eqref{eq:chi_gap},
$|\chi_m-\chi_n|\ge\kappa_1\max(1,c)|m-n|$,
valid uniformly across Regimes~(A) and~(B) with $\kappa_1=1/\pi$.

\noindent\textbf{Assembly.}
\[
  |\Phi_{mn}'(x)|
  =\frac{|\chi_m-\chi_n|}{T(\sqrt\mu+\sqrt\nu)}
  \ge\frac{\kappa_1 c|m-n|}{T\cdot C_\nu c}
  =\frac{\kappa_1}{C_\nu T}|m-n|
  =:\frac{\kappa}{T}|m-n|.
\]
The factor $c$ cancels exactly; $\kappa=\kappa_1/C_\nu$ is
independent of $c$ and of which regime applies.

%% FIX 3: non-sharp denominator declaration
The upper bound $\sqrt{\mu}+\sqrt{\nu}\le C_\nu c$ uses
$\chi_k\le C_{\gamma_0}c^2$ for $k\in\mathcal{B}$, which is coarse
for $n\ll c$ (where $\sqrt{\chi_n}\sim n\ll c$). This
non-sharpness is intentional: a tighter bound would yield a larger
$\kappa$ but is not needed, since only positivity of $\kappa$ is
used in the IBP argument (Lemma~\ref{lem:ibp_uniform}).
The IBP argument depends only on positivity of $\kappa$, not on
its precise value.
\end{proof}

\begin{remark}[Bulk restriction]
Lemma~\ref{lem:phase_gap_final} holds only for
$m,n\in\mathcal{B}$. In the transition regime
$|k-N_{\mathrm{Sh}}|=O(c^{1/3})$ the spectral gap degrades to
$O(c^{1/3})$ and the argument fails; that regime is excluded
throughout Section~\ref{sec:decay}.
\end{remark}

\begin{lemma}[$c$-uniform amplitude derivative control]%
\label{lem:amplitude_control}
The Pr\"ufer amplitude $r_n(x)$ of $\psi_n^{(c)}$ satisfies, for
$x\in\Rbulk^{(m,n)}$ and all $k\ge 0$:
\[
  \|\partial_x^k r_n\|_{L^\infty(\Rbulk^{(m,n)})}\le C_k,
\]
with $C_k$ depending on $k$, $\rho$, $\gamma_0$, $\delta$,
but \emph{not} on $c\ge c_0$ or $n\in\mathcal{B}$.
The WKB amplitude product $a_{mn}:=r_m\overline{r_n}$ satisfies
$\|\partial_x^k a_{mn}\|_{L^\infty(\Rbulk^{(m,n)})}\le C_k'$ with the
same uniformity.
\end{lemma}
\begin{proof}
By Lemma~\ref{lem:discrete_continuous_bridge}(iv),
$r_n(x)=r_n^{\mathrm{WKB}}(x)(1+O(1/n))$ uniformly on
$I_\delta\supset\Rbulk^{(m,n)}$. The WKB amplitude is
$r_n^{\mathrm{WKB}}(x)\propto(\omega_n^{\mathrm{cl}}(x))^{-1/2}$
\cite[Definition~4.1]{PaperIV}; its $k$-th derivative involves only
$(\omega_n^{\mathrm{cl}})^{(j)}/(\omega_n^{\mathrm{cl}})^{j/2+1/2}$
for $j\le k$. By Lemma~\ref{lem:discrete_continuous_bridge}(i),
$\omega_n^{\mathrm{cl}}\asymp n$ and all its spatial derivatives are
bounded by $C/n^j$ on $I_\delta$ uniformly in $c$
\cite[Lemma~2.4]{PaperIV}. Consequently, each derivative of
$r_n^{\mathrm{WKB}}$ is $O(1)$ in $c$: the $n$-factors cancel against
the normalisation, and no $c^{k/2}$ growth occurs. The claim for
$a_{mn}=r_m\overline{r_n}$ follows by the product rule.
\end{proof}

\begin{remark}[Why no $c^{k/2}$ appears]
Lemma~\ref{lem:amplitude_control} separates the oscillatory factor
$e^{i\theta_n}$ from the amplitude $r_n$. All $c$-dependence resides
in the phase $\theta_n'(x)\asymp c/T$; once this factor is absorbed
by the IBP operator $L=(i\Phi_{mn}')^{-1}\partial_x$, the remaining
amplitude chain is $c$-free. The erroneous appearance of $c^{k/2}$
in earlier drafts arose from differentiating
$\psi_n^{(c)}=r_n\sin\theta_n$ without first separating phase from
amplitude.
\end{remark}

\begin{lemma}[Bulk-localised IBP control]\label{lem:ibp_uniform}
Let $L:=(i\Phi_{mn}')^{-1}\partial_x$ on $\Rbulk^{(m,n)}$.
For every $p\ge 1$:
\[
  \|(L^*)^p a_{mn}\|_{L^1(\Rbulk^{(m,n)})}
  \le C_p\left(\frac{T}{|m-n|}\right)^p,
\]
with constants independent of $m,n\in\mathcal{B}$, $m\ne n$,
$c\ge c_0$, and $p\le p_0$.
\end{lemma}
\begin{proof}
%% FIX 4: explicit regularity of Phi' and well-definedness of L
Since $\Phi_{mn}'(x)=(\chi_m-\chi_n)/[T(\sqrt{\mu}+\sqrt{\nu})]$
and $\sqrt{\mu}+\sqrt{\nu}$ is bounded away from zero on
$\Rbulk^{(m,n)}$ by Lemma~\ref{lem:bulk_spectral_sep}, the function
$\Phi_{mn}'$ is $C^\infty$ on $\Rbulk^{(m,n)}$ and
$(\Phi_{mn}')^{-1}$ is bounded; hence $L$ and its adjoint $L^*$
are well-defined bounded operators on $C^\infty(\Rbulk^{(m,n)})$.

\noindent\textbf{Structure.}
By induction, $(L^*)^p a_{mn}$ is a linear combination of
$\partial_x^k a_{mn}/(\Phi_{mn}')^\ell$, $k\le p$, $\ell\le p$.
Each such term is bounded by
$C_k'\cdot(T/(\kappa|m-n|))^\ell\le C_p(T/|m-n|)^p$,
using Lemma~\ref{lem:amplitude_control} for the numerator
($c$-independent) and Lemma~\ref{lem:phase_gap_final} for the
denominator ($\kappa$ uniform across both regimes).
Integration over $\Rbulk^{(m,n)}\subset[-T,T]$
(length $\le 2T$) gives the $L^1$ bound.

\noindent\textbf{Boundary terms.}
At $x\in\partial\Rbulk^{(m,n)}$, the Olver matching gives
$|a_{mn}(x)|\le C_\infty^2$ and
$|(\Phi_{mn}')^{-1}|\le T/(\kappa|m-n|)$,
so the boundary contribution is $O(T/|m-n|)^1$ and is absorbed into
the $p=1$ bound.
\end{proof}

\begin{lemma}[Turning-point estimate]\label{lem:turning_revised}
For all $0\le m,n\le N-1$:
$\int_{U_m}|\psi_m^{(c)}\psi_n^{(c)}|\,dx\le C\,c^{-1/3}$.
\end{lemma}
\begin{proof}
$\|\psi_k^{(c)}\|_{L^\infty}\le C_\infty$ \cite{BonamiKaroui2014}
and $|U_m|=O(c^{-1/3})$; Cauchy--Schwarz gives $O(c^{-1/3})$.
\end{proof}

\begin{lemma}[Bulk decay for the continuous integral]%
\label{lem:bulk_decay}
For $m\ne n$ in $\mathcal{B}$ and every $p\ge 1$:
\[
  \left|\int_{\Rbulk^{(m,n)}}f_{mn}(x)\,dx\right|
  \le C_p\left(\frac{T}{|m-n|}\right)^p.
\]
The bound is $c$-independent. No transition to
$(1+|m-n|)^{-p}$ is made; that step belongs to DSTP.
\end{lemma}
\begin{proof}
Apply $L$ iteratively $p$ times to
$\int_{\Rbulk^{(m,n)}}a_{mn}(x)e^{i\Phi_{mn}(x)}\,dx$,
using Lemma~\ref{lem:ibp_uniform}.
\end{proof}

\begin{assumption}[First-order quadrature consistency]%
\label{ass:quad_first_order}
There exists $C_Q>0$ such that for all $f\in W^{1,1}([-T,T])$:
\[
  \Bigl|\int_{-T}^T f\,dx-\sum_j w_j f(p_j)\Bigr|
  \le C_Q h\int_{-T}^T|f'|\,dx,
  \qquad h:=\max_j|p_{j+1}-p_j|.
\]
\end{assumption}

\begin{lemma}[Derivative bound]\label{lem:derivative_bound}
$\int_{-T}^T|f'_{mn}|\,dx\le Cc$, uniformly in $m,n,c$.
\end{lemma}
\begin{proof}
$|f'_{mn}|\le|\psi_m'||\psi_n|+|\psi_m||\psi_n'|\le 2C_0C_1 c$
by \cite{BonamiKaroui2014}.
\end{proof}

\begin{theorem}[Quadrature defect bound]\label{thm:banded_decay}
Under Assumption~\ref{ass:quad_first_order},
$|E_{mn}|\le Chc$ for all $0\le m,n\le N-1$.
\end{theorem}
\begin{proof}
$|E_{mn}|=|R_{mn}^{\mathrm{quad}}|
\le C_Q h\int_{-T}^T|f'_{mn}|\,dx\le C_QChc$.
\end{proof}

\begin{lemma}[Continuous bulk--turning decay]%
\label{lem:bulk_decay_continuous}
For $m\ne n$ in $\mathcal{B}$ and every $p\ge 1$:
\[
  \left|\int_{-T}^T\psi_m^{(c)}\overline{\psi_n^{(c)}}\,dx\right|
  \le C_p\left(\frac{T}{|m-n|}\right)^p+C\,c^{-1/3}.
\]
\end{lemma}
\begin{proof}
Decompose into $\Rbulk^{(m,n)}$ and $\Rturning^{(m,n)}$; apply
Lemmas~\ref{lem:bulk_decay} and~\ref{lem:turning_revised}.
\end{proof}

\begin{remark}[Separation of continuous and discrete decay]
The bound of Lemma~\ref{lem:bulk_decay_continuous} describes
the \emph{continuous} PSWF kernel. It does not transfer to
$E_{mn}=R_{mn}^{\mathrm{quad}}$ without additional quadrature
structure: generic rules destroy the oscillatory cancellation,
yielding only the structure-blind $|E_{mn}|\le Chc$. The gap is
precisely what DSTP (Definition~\ref{def:dstp}) encodes.
\end{remark}

%% ==================================================================
\section{Frame stability}\label{sec:stability}
%% ==================================================================

\begin{theorem}[Defect decomposition]\label{thm:main_reduction}
\[
  \EN=M^*(D-I)M+(G^{\mathrm{w}}_{N,c}(\mathbf{p},I)
     -G^{\mathrm{ref}}_{N,c}),
\]
\[
  \norm{\EN}_2\le\norm{D-I}_2\cdot N\omega/\pi
  +\norm{G^{\mathrm{w}}_{N,c}(\mathbf{p},I)
         -G^{\mathrm{ref}}_{N,c}}_2.
\]
\end{theorem}
\begin{proof}
Add and subtract $M^*M$; apply Lemma~\ref{lem:lambda_max_bound}.
\end{proof}

\begin{lemma}[Weyl perturbation]\label{lem:weyl}
If $\norm{\EN}_2<\lambda_{N-1}(c)$ then
$\lambda_{\min}(G^{\mathrm{w}}_{N,c})
\ge\lambda_{N-1}(c)-\norm{\EN}_2>0$.
\end{lemma}
\begin{proof}
Weyl's inequality: eigenvalues of $G^{\mathrm{ref}}_{N,c}+\EN$
interlace those of $G^{\mathrm{ref}}_{N,c}$ up to $\norm{\EN}_2$.
\end{proof}

\begin{definition}[DSTP]\label{def:dstp}
The rule $\{(p_j,w_j)\}$ satisfies the \emph{Discrete Spectral
Transfer Property} on $V_{N,c}$ if there exist $C>0$, $p>1$ such that
\[
  |E_{mn}|\le Chc\,(1+|m-n|)^{-p}
  \qquad\forall\,0\le m,n\le N-1.
\]
DSTP is an axiomatic interface hypothesis, strictly beyond
Assumption~\ref{ass:quad_first_order} and not implied by
Lemma~\ref{lem:bulk_decay_continuous}. It encodes the missing
oscillatory transfer mechanism. Verification for concrete schemes
is deferred to~\cite{PaperIIquad}.
\end{definition}

\begin{remark}[DSTP heuristics]
DSTP is not expected for Riemann sums, Newton--Cotes, or unadapted
Gauss--Legendre rules, which carry no phase-alignment with the PSWF
eigenbasis. It is a structural strengthening of MZ-type conditions,
analogous in spirit to spectral exactness requirements, specialised
to the off-diagonal decay geometry of PSWF kernels.
\end{remark}

\begin{remark}[Non-reducibility of \textup{DSTP} --- constructive witness]
\label{rem:ax5_independence}
While \textup{AX1--AX4} control the \emph{analytic} structure of the
system (symbol class, spectral gaps, quadrature order, uniform PSWF
bounds), DSTP controls the \emph{geometric} compatibility of the
discretisation. The two layers are logically independent within this
framework, as the following constructive witness shows.

\smallskip
\noindent\textbf{Witness rule.}
Let $\{p_j=-T+2Tj/N\}_{j=1}^N$ be the equidistant Riemann rule
with weights $w_j=2T/N$ and mesh $h=2T/N$.
This rule satisfies:
\begin{enumerate}[(i)]
\item \textup{AX3} (Assumption~\ref{ass:quad_first_order}),
  with $C_Q=1$;
\item \textup{AX4} (uniform PSWF bounds \cite{BonamiKaroui2014}),
  since $\|\psi_n^{(c)}\|_{L^\infty}$ and
  $\|(\psi_n^{(c)})'\|_{L^\infty}$ are independent of node locations;
\item \textup{AX1--AX2} (symbol class and spectral gap), from
  \cite{PaperIV,OsipovRokhlinXiao} alone.
\end{enumerate}

\smallskip
\noindent\textbf{Spectral concentration of $f_{mn}$.}
Since $\hat{f}_{mn}=\hat{\psi}_m^{(c)}\ast\hat{\psi}_n^{(c)}$
and each $\hat{\psi}_k^{(c)}$ is exponentially concentrated on
$[-\omega/\pi,\omega/\pi]$ \cite{BonamiKaroui}, the convolution
satisfies
\[
  \int_{|\xi|>4c/T}|\hat{f}_{mn}(\xi)|^2\,d\xi\le Ce^{-\alpha c},
\]
with $\alpha,C>0$ independent of $m,n,c$.
In particular $\hat{f}_{mn}(k/h)=O(e^{-\alpha c})$ for $|k|\ge 2$,
so the aliasing sum is dominated by the $k=1$ term.

\smallskip
\noindent\textbf{DSTP fails.}
By the Poisson summation formula,
$h\sum_j f_{mn}(p_j)=\int_{-T}^T f_{mn}\,dx
+\sum_{k\ne 0}\hat{f}_{mn}(k/h)$.
Since $1/h\sim c/(2T)$ lies within the spectral concentration band
$[0,4c/T]$, the $k=1$ term is not exponentially small.
Accumulating over off-diagonal pairs:
\[
  \sum_{m\ne n}|E_{mn}^{\mathrm{Riemann}}|^2
  \ge\sum_{m\ne n}|\hat{f}_{mn}(1/h)|^2-O(N^2e^{-\alpha c})
  \ge c_0(hc)^2N,
\]
therefore
\[
  \|E_{N,c}^{\mathrm{Riemann}}\|_2
  \ge\frac{\|E_{N,c}^{\mathrm{Riemann}}\|_F}{\sqrt{N}}
  \ge c_0^{1/2}hc,
\]
with no decay in $|m-n|$, so DSTP fails for all $p>0$.

\smallskip
\noindent\textbf{Structural interpretation.}
\textup{AX1--AX4} supply $c$-uniform amplitude control, spectral gaps,
quadrature consistency, and PSWF bounds --- yet none encodes geometric
compatibility between $\{p_j\}$ and the oscillatory eigenstructure
of $V_{N,c}$.
\begin{quotation}
  \noindent\textit{DSTP is not an estimate. It is a spectral
  incompatibility condition: the aliasing frequency of the sampling
  rule falls within the spectral concentration band of the PSWF
  product kernel.}
\end{quotation}
This places DSTP in the same logical category as spectral exactness
in Gauss quadrature, sampling compatibility in the Shannon--Nyquist
theorem, and frame reconstruction axioms in non-uniform sampling
theory --- none of which follow from regularity bounds alone.

\smallskip
\noindent\textbf{Minimality.}
DSTP is \emph{orthogonal} to \textup{AX1--AX4} in the sense of
logical independence of admissible discretisation classes within
this framework. The Riemann-rule witness shows a rule may satisfy
all of \textup{AX1--AX4} and fail DSTP; no combination of analytic
bounds implies geometric compatibility.
The axiom system $\{\textup{AX1},\dots,\textup{AX4},\textup{DSTP}\}$
is therefore minimal: no axiom is derivable from the remaining ones
within this framework.

Verification of DSTP for PSWF-adapted (XRY/Gauss--PSWF) quadrature
schemes is the content of~\cite{PaperIIquad}.
\end{remark}

\begin{theorem}[Conditional frame stability]%
\label{thm:conditional_frame_schur}
Let Assumption~\ref{ass:quad_first_order} and DSTP hold with $N\sim c$,
$p>1$. If $Chc\,\zeta(p)<\lambda_{N-1}(c)$,
where $\zeta(p):=\sum_{k\in\mathbb{Z}}(1+|k|)^{-p}<\infty$, then
$\{(p_j,w_j)\}$ is a weighted frame for $V_{N,c}$ with bounds
\begin{align*}
  A(N,c)&:=\lambda_{N-1}(c)-Chc\,\zeta(p)>0,\\
  B(N,c)&:=\lambda_0(c)+Chc\,\zeta(p).
\end{align*}
This result is optimal conditional on DSTP: no stronger conclusion
is possible from first-order quadrature consistency alone.
\end{theorem}
\begin{proof}
DSTP and the Schur test give
$\norm{\EN}_2\le Chc\,\zeta(p)$.
Apply Lemma~\ref{lem:weyl}.
\end{proof}

\begin{remark}[Asymptotic scope]\label{rem:asymptotic_scope}
The frame condition $Chc\,\zeta(p)<\lambda_{N-1}(c)$ requires
$h\ll\lambda_{N-1}(c)/c$. Since $\lambda_{N-1}(c)\sim e^{-\alpha c}$
for $N\sim c$ \cite{BonamiKaroui}, this demands exponentially small
mesh size in the transition regime. This reflects the intrinsic
difficulty of the PSWF spectral edge and is recorded as the
fundamental quantitative obstruction.
\end{remark}

\begin{corollary}\label{cor:fixed_param_scope}
Theorem~\ref{thm:conditional_frame_schur} yields finite-dimensional
frame stability under DSTP. The structure-blind bound
$|E_{mn}|\le Chc$ alone does not imply Schur-summable decay.
\end{corollary}

%% ==================================================================
\section{Boundary to the companion framework}\label{sec:bridge}
%% ==================================================================

This paper proves a reduction principle: discrete frame stability is
reduced to the single interface hypothesis DSTP, with the obstruction
made explicit. The continuous bulk--turning analysis of
Section~\ref{sec:decay} identifies the oscillatory decay mechanism
and demonstrates that it is a purely continuous phenomenon destroyed
by generic quadrature. It is not used in the proof of
Theorem~\ref{thm:conditional_frame_schur}; its role is to motivate
DSTP as the necessary interface axiom.

Three mechanisms are cleanly separated:
(1)~continuous PSWF orthogonality (exact,
Proposition~\ref{prop:xry_split});
(2)~oscillatory decay of the continuous kernel
(Lemma~\ref{lem:bulk_decay_continuous}, bulk sector only);
(3)~quadrature-induced defect (structure-blind $|E_{mn}|\le Chc$,
global). The companion paper~\cite{PaperIIquad} addresses whether
DSTP holds for PSWF-adapted quadrature schemes.

%% ==================================================================
\section{Open problems}\label{sec:open}
%% ==================================================================

\begin{problem}[Exponential decay in bulk]
Upgrade the algebraic bound $(T/|m-n|)^p$ to exponential decay
$e^{-\alpha|m-n|}$ in the strict bulk; see Paper~V \cite{PaperV}.
\end{problem}

\begin{problem}[Sharper turning bounds]
Refine $O(c^{-1/3})$ via microlocal analysis near turning points.
\end{problem}

\begin{problem}[Explicit constants]
Make all constants in
Lemmas~\ref{lem:bulk_spectral_sep}--\ref{lem:ibp_uniform}
fully explicit in $N$ and $c$.
\end{problem}

\begin{problem}[Oscillation-compatible quadrature]
Construct concrete schemes satisfying DSTP; see \cite{PaperIIquad}.
\end{problem}

\begin{problem}[Uniform-grid unisolvence]
Determine whether uniform grids are unisolvent for $V_{N,c}$.
\end{problem}

\begin{problem}[Coupled limits]
Analyse $\lambda_{\min}(G^{\mathrm{w}}_{N,c})$ for $N,c\to\infty$
simultaneously.
\end{problem}

\begin{problem}[Near-base regime]
Sharpen Proposition~\ref{prop:spectral_gap} Regime~(B)
($\min(m,n)<c^{1/2}$) to obtain uniform $c$-dependence at the
spectral base; the current bound $|\chi_m-\chi_n|\ge|m-n|$ is
$c$-independent but not $c$-enhanced.
\end{problem}

\begin{problem}[Universal necessity of DSTP]
Determine whether every quadrature rule yielding a stable frame for
$V_{N,c}$ (in the sense of Definition~\ref{def:framestable}) must
satisfy DSTP or an equivalent spectral compatibility condition.
The constructive witness of Remark~\ref{rem:ax5_independence} shows
sufficiency of \textup{AX1--AX4} is not enough; a converse
characterisation remains open.
\end{problem}

%% ==================================================================
\appendix
\section{Airy uniformisation}\label{app:airy}
%% ==================================================================

\begin{lemma}[PSWF Airy approximation]\label{lem:airy_approx}
For $0\le n\le N-1$ and $x\in U_n$:
\[
  \psi_n^{(c)}(x)
  =\kappa_n^{(c)}\,\mathrm{Ai}(-\zeta_n(x))+O(c^{-1/3}),
\]
$\zeta_n(x)=c^{2/3}(x-x_n^*)/T+O(c^{-1/3})$,
$|\kappa_n^{(c)}|\le C_\kappa$ uniformly
\cite{Olver1974,OsipovRokhlinXiao}.
\end{lemma}

\begin{lemma}[Turning-point integral via Airy]\label{lem:airy_bound}
$|\int_{U_m\cup U_n}\psi_m^{(c)}\overline{\psi_n^{(c)}}\,dx|
\le C\,c^{-1/3}$.
\end{lemma}
\begin{proof}
$\|\psi_k^{(c)}\|_{L^2(U_m\cup U_n)}=O(c^{-1/6})$; Cauchy--Schwarz.
\end{proof}

%% ==================================================================
\begin{thebibliography}{99}

\bibitem{BonamiKaroui}
A.~Bonami and A.~Karoui,
\textit{Spectral decay of the sinc kernel operator and approximation
with prolate spheroidal wave functions}, arXiv:1012.3881.

\bibitem{BonamiKaroui2014}
A.~Bonami and A.~Karoui,
\textit{Uniform approximation and explicit estimates for the prolate
spheroidal wave functions}, arXiv:1405.3676.

\bibitem{Olver1974}
F.~W.~J.~Olver,
\textit{Asymptotics and Special Functions}, Academic Press, 1974.

\bibitem{OsipovRokhlinXiao}
A.~Osipov, V.~Rokhlin, and H.~Xiao,
\textit{Prolate Spheroidal Wave Functions of Order Zero},
Springer, 2013.

\bibitem{PaperII}
Anonymous Author,
\emph{Scaling Limits, Trace Formula, and Weil-Operator Connections
for Prolate Sampling Operators}, preprint, 2026.

\bibitem{PaperIIquad}
Anonymous Author,
\emph{Conditional Frame Stability via PSWF Spectral Projection
and XRY Quadrature}, preprint, 2026.

\bibitem{PaperIV}
Anonymous Author,
\emph{Semiclassical Equidistribution of PSWF Densities via
Pr\"{u}fer Phase Analysis}, preprint, 2026.

\bibitem{PaperV}
Anonymous Author,
\emph{WKB Stability and Exponential Concentration in the
Prolate Bulk Sector}, preprint, 2026.

\bibitem{Slepian1978}
D.~Slepian,
\textit{Prolate spheroidal wave functions, Fourier analysis,
and uncertainty~-- V},
Bell Syst.\ Tech.\ J.\ \textbf{57} (1978), 1371--1430.

\bibitem{Xiao2001}
H.~Xiao, V.~Rokhlin, and N.~Yarvin,
\textit{Prolate spheroidal wave functions, quadrature, and
interpolation},
Inverse Problems \textbf{17} (2001), 805--838.

\end{thebibliography}

\end{document}
